% EGA V, English translation by P. Blass, 1987 (cote 24): archive cote 24, site page 8 (PDF page 9).
% typescript. Draft compared: none.
% AI-transcribed by Claude Opus 5.5 (claude-opus-5-5, adaptive thinking, effort high), Message Batches API, batch ega24b (msgbatch_01RGJJeHdR96U16ZWFuUjtDx), 2026-11-13; not yet checked by eye.
\scanpage{24/8}
\mspage{8}

(the squares and diamonds appearing in this diagram are
Cartesian). In the next section we will study systematically the
following case: $S^{1}$ is the spectrum of a field K and its image in P
is generic in the corresponding fiber $P^{V}_{\unsure{S|s}}$. After making the base
change Spec k(s) $\rightarrow$ S we are reduced to the case where S is the
spectrum of a field k, which is what we assume in the next
section. Also the majority of properties studied for X and Y are
of ``geometric nature'' and therefore invariant under base change,
which allows us also (without loss of generality) to limit
ourselves to the case where K is algebraically closed or to the
case where K = k($\eta$), $\eta$ , being the generic point of $P^{V}$ and . $\xi$ :
Spec(K) $\rightarrow$ $P^{V}$ is of course the canonical morphism. We also note
that for geometric questions concerning X, Y we can (after making
a base change) restrict ourselves to the case of k algebraically
closed. A terminological reminder. If f is an immersion we
usually call Y a hyperplane section of X (relative to the
projective immersion f and the hyperplane H). There is no reason
not to extend this terminology to the case of an arbitrary f.
